\chapter{Limitations}
\label{ch:lim}

% Written 21 September 2026 from analysis/SPINE.md sections 6, 7b, 8 and 10.

Each limitation below is stated as the threat, the claim it weakens, and
what would remove it.
Limitations that the results themselves impose are marked as such; the
others were fixed by the design and were known before the runs.

\section{Scope}
\label{sec:lim-scope}

One base model (Gemma-2-2b-it), one adapter rank, one environment family,
100 epochs, and three to five seeds per arm.
The hyperparameters are this stack's and differ from ShapeLLM's released
configuration in GAE \(\lambda\), value coefficient, mini-batch, score
scaling, entropy and the shaper's learning rate and clip
(Chapter~\ref{ch:design}).
\textbf{Weakens:} any statement of the form ``opponent shaping fails for
language-model agents''.
Nothing here licenses that sentence; the result is about one objective, one
estimator of it, and one environment.
\textbf{Removed by:} the same design at a second model scale and a second
environment family.

\section{The control already succeeded}
\label{sec:lim-headroom}

At \(m=2\) two naive learners produced an unconditional co-learner on five
of five seeds (Chapter~\ref{ch:results}).
The solver's stake --- a restrained partner is worth 55.1 units to a best
responder --- measures what restraint is worth, not how much room a shaper
had to add restraint that was not already there.
The room was smaller than the stake suggests, and condition~2 of S failed on
the controls rather than on the shapers.
\textbf{Weakens:} the design's power to detect a positive effect, and hence
the strength of the null as evidence that the objective does nothing.
\textbf{Removed by:} a prior, budget or environment in which the naive
baseline does not acquire the target behaviour on its own --- at the cost of
a new pre-registration, because changing the design point after seeing these
results would convert a registered null into an unregistered positive.

\section{Objective and estimator are confounded in the pre-registered arms}
\label{sec:lim-estimator}

Every pre-registered shaper arm estimates cross-episode credit the same
way, by chaining GAE across the trial.
The split-credit arm, which does not, is an addition made after the second
gate failure and is labelled as such.
The experiment therefore compares one objective under one estimator against
its controls, and separately reports a second estimator whose inclusion was
not pre-registered.
\textbf{Weakens:} reading the null as a property of the trial-level
objective rather than of chained whole-trial GAE.
\textbf{Removed by:} a pre-registered comparison of two or more estimators
of the same objective, with the estimator as the manipulated factor.

\section{The implemented objective is not the ideal one}
\label{sec:lim-attenuation}

Two implementation facts bound what the cross-episode term can carry.
Credit reaches a step discounted by \(\lambda\) per live step, so a reward
one episode later enters the opening advantage of a full episode at
\(0.97^{50}=0.22\), and later episodes arrive more weakly still.
Advantages are whitened per update, which subtracts the trial's mean and so
removes the component of the term that is common to the whole trial ---
under equal episode lengths and a persistent co-learner response, roughly
three fifths of that response's expected contribution.
\textbf{Weakens:} the identification of ``the trial-level objective'' with
what was optimised.
What was optimised is a \(\lambda\)-attenuated, trial-mean-centred proxy for
it.
\textbf{Removed by:} adding the cross-episode term after normalisation, or
batching so that the trial-common component survives, with the objective
otherwise unchanged.

\section{The probe cannot classify at \texorpdfstring{\(m=3\)}{m = 3}}
\label{sec:lim-m3-probe}

The probe splits rounds at stock 12.
At \(m=3\) a restrained pool sits at capacity, so between 0.004 and 0.3 per
cent of live rounds fall below that cut for the slow and ShapeLLM-style
arms: \(r_{\mathrm{low}}\) rests on 1, 8, 31 or 72 rounds, and two \(m=3\)
probes are undetermined (Chapter~\ref{ch:results}).
\textbf{Weakens:} every class label at \(m=3\), including the two
``unconditional'' labels.
It does not weaken decision~C, which reads \(r_{\mathrm{after\ take}}\) on
between 10{,}250 and 12{,}340 rounds per probe, nor the behavioural
readouts.
\textbf{Removed by:} a regime-specific stock cut fixed in advance, or a
probe partner that draws the stock down so both bands are populated.

\section{Seeds and the decision rule}
\label{sec:lim-seeds}

``Most seeds'' is three of five or two of three, applied to four
simultaneous conditions with no test statistic.
The rule is deliberately conservative: it is hard to pass by chance and easy
to miss a partial effect.
Two vetoes --- the slow arm at \(m=2\) and decision~C at \(m=3\) --- rest on
three seeds, where a single seed decides a class count.
The Wilson intervals reported with each rate are within-run and carry no
run-to-run variance.
\textbf{Weakens:} the null as evidence of absence rather than absence of
evidence, and the two three-seed vetoes.
\textbf{Removed by:} five seeds on every arm, declared before the runs so
that seeds are not added until a count flips, with the analysis rerun on the
full set.

\section{Transfer has no control yet}
\label{sec:lim-e1}

Condition~3 of S reads the class of a fresh learner trained against the
frozen shaper, but E1 exists only for shapers.
It cannot presently separate \emph{this shaper's policy teaches a stranger}
from \emph{any trained agent~2 teaches a stranger}, and the executed E2 tape
result makes the second reading live.
\textbf{Weakens:} condition~3, and with it any transfer claim.
\textbf{Removed by:} the control runs specified in the pre-registration's
amendment of 21~September --- a fresh learner against the frozen
trial-batched agent~2, five seeds --- which are under way.

\section{The dial is exploratory}
\label{sec:lim-exploratory}

The design point was amended after a first gate failure, the estimator
repair that produced the split-credit arm was chosen after a second, and
training followed a gate that passed at the third attempt.
Table~\ref{tab:evidence-status} in Chapter~\ref{ch:design} states which parts were
fixed in advance.

\textbf{Weakens:} the confirmatory status of everything downstream of the
second gate failure.
\textbf{Removed by:} a replication of the executed design at a new
pre-registration, which is the honest next step rather than a reanalysis of
these runs.

\section{The failure is located but not isolated}
\label{sec:lim-diagnostic}

The evidence that the estimator rather than the game produced the collapse
is indirect: the gates show the environment is learnable, the channel
measure shows no response specific to the objective, and the split-credit
arm shows that changing only the estimator changes the outcome.
No run isolates the chained estimator from the co-learning dynamics.
\textbf{Weakens:} the mechanism argument of Chapter~\ref{ch:discussion}.
\textbf{Removed by:} one cheap diagnostic --- the chained-credit shaper
against the scripted tit-for-tat partner of gate G2, three seeds, 100
epochs, about four GPU-hours.
A naive learner acquires reciprocity against that partner
(Section~\ref{sec:gates}); if the chained shaper does not, the failure is in
the estimator and not in the presence of a learning co-player.
The run has no co-learner and therefore cannot produce a shaping result; it
can only locate the failure.
